% =====================================================================
%  Harald Høffding, Review of H. Martensen, Jacob Böhme
%  Translated from the Danish: „H. Martensen: Jacob Böhme (Anmeldelse)“, Ude og Hjemme, vol. 5 (1881--82), no. 238 [23 April 1882], pp. 365--366.
%
%  TRANSLATION -- GENERATED; DO NOT EDIT BY HAND.  Built by ebook-method/hoeffding-1879-82/tr_parts.py build boehme-anmeldelse-1882 --write
%  from the parts in ebook-method/hoeffding-1879-82/tr/ (GLOSSARY.md, TRANSLATOR-BRIEF.md).  Made from the collated
%  transcription texts/hoeffding/boehme-anmeldelse-1882/transcription.tex (2026-10-15); [N] in the margin = the page of the original.
%  Terms: GLOSSARY.md (choices marked ◆ there for Hans to review).
% =====================================================================
\documentclass[12pt,a4paper,openany]{book}
\usepackage[utf8]{inputenc}
\usepackage[T1]{fontenc}
\usepackage{libertinus}
\usepackage{libertinust1math}
\usepackage[german,french,english]{babel}
\usepackage{textalpha}
\usepackage[protrusion=true,expansion=true]{microtype}
\usepackage{setspace}
\usepackage[margin=1.3in]{geometry}
\usepackage{fancyhdr}
\usepackage{hyperref}
\hypersetup{pdftitle={Review of Martensen, Jacob Böhme}, pdfauthor={Harald Høffding}, hidelinks}
\pagestyle{fancy}
\fancyhf{}
\fancyhead[LE,RO]{\thepage}
\fancyhead[C]{\textit{H. Høffding: Review of Martensen, Jacob Böhme (1882)}}
\fancypagestyle{plain}{\fancyhf{}\fancyhead[LE,RO]{\thepage}}
\setstretch{1.3}
\usepackage{marginnote}
\usepackage{graphicx}
\DeclareUnicodeCharacter{0254}{\reflectbox{c}}   % ɔ (id est)

\newcommand{\opage}[1]{\marginnote{\footnotesize\textit{[#1]}}}
\newcommand{\signature}[1]{\par\smallskip\hfill #1\hspace*{2em}\par}
\pdfstringdefDisableCommands{\renewcommand{\footnote}[2][]{}\renewcommand{\opage}[1]{}}

\begin{document}

\begin{titlepage}
\begin{center}
{\Huge \emph{Ude og Hjemme}}\\[0.4em]{\large [Abroad and at Home]}\\[0.6em]{\large A Nordic Illustrated Weekly}\\[0.6em]{\small edited by Otto Borchsenius and F. Hendriksen}\\[0.6em]{\small Fifth Volume, 1881--1882}\\[0.4em]{\small Copenhagen: F. Hendriksen}\\[2em]
{\small Translated from the Danish}
\end{center}
\end{titlepage}
\clearpage

% --- p. 365 ---
\setcounter{footnote}{0}%
\opage{365}\begin{center}{\large H. \emph{Martensen: Jacob Böhme.}}\\[0.3em]{\small Theosofiske Studier [Theosophical Studies].  Copenhagen: Gyldendalske Boghandels Forlag, 1881.}\end{center}\medskip

% (p. 365, column 1)
\setcounter{footnote}{0}%
The founding of the Protestant Church was one of the most significant results of the general yearning for liberation, inwardness and naturalness that distinguished the age that has been called the age of rebirth (the Renaissance). But scarcely had the young Church secured its existence before it thought it possessed freedom and inwardness enough. It formulated its dogmas and pushed back every free direction of thought and feeling. Although Luther himself was deeply influenced by the mystical-speculative current that had unfolded in Germany from the 14th century on, and whose most outstanding thinker was Master Eckhart, the mystics who arose within the new Church nevertheless suffered rough treatment. Schwenckfeld and Sebastian Franck led a wandering and unsettled life for the sake of their views, and Valentin Weigel lived undisturbed only because he did not let his writings be printed. At the universities orthodoxy ruled, and the external combination of theology and Aristotelian philosophy founded by Melanchton. But the mystical books and manuscripts passed from hand to hand among laymen. Probably the remarkable shoemaker Jacob Böhme of Görlitz (b. 1575, d. 1624) was initiated into these writings while he traveled on his trade. Journeymen artisans on the road were among the most important bearers and disseminators of the mystical ideas and writings. With religious mysticism there already combined in Weigel the influence of Paracelsus' alchemy and philosophy of nature, as a realistic ferment that was added to the mystical idealism. This double movement was carried to its highest and most significant form when it struck down into Böhme's deep, childlike mind. Religious brooding, grand imagination, lively sense of nature, psychological insight and keen thinking, in union with a marvelous originality and power in language, made him a unique figure in the history of thought. He is the mythological predecessor of German philosophy, whose fundamental ideas he utters in mystic form.

% (p. 365, column 2)
\setcounter{footnote}{0}%
It is this man whose life and doctrine form the main content of the present work. There is a mental kinship between the author and his hero. Martensen's thinking has always had a decidedly mystical or theosophical stamp. When in theology the word so often must serve in place of the missing concept, then with Martensen it is surely most often the pregnant image, the imagination full of presentiment, that forms the conclusion of the thought and veils what it is unable to fill out. What places him so high above the whole younger generation of theologians is that he nevertheless still has the courage to think. He reproaches the younger theologians for merely keeping to the received doctrine of the Church and not engaging in winning a coherent knowledge and understanding. It is beautiful to see the old champion of theological speculation remain true to the love of his youth to the last. Just as one of the works with which he began his brilliant authorship was a work on Master Eckhart, so he has now in his high old age seated himself at the feet of the profound shoemaker from Görlitz to depict the concluding representative of that mystical-theosophical movement and himself to spin the thread further. For the work is not a purely historical monograph; the author does not merely criticize his hero, but seeks also to go further along the tracks pointed out by him.

Where the Bishop and the Shoemaker disagree, I for my part am most inclined to take the Shoemaker's side; and where they agree, I often disagree with them both. But it does seem to me that one reads the present work with far greater interest and profit than the author's ``Christian Ethics.'' In the ``Ethics'' one moved in the real world and saw the sharpened high-church line of thought collide with all the strivings for freedom of the present day. Here, by contrast, one moves so high above reality that only on rare occasions one catches sight of the earth as a vanishing atom. We are led into
% --- p. 366 ---
\setcounter{footnote}{0}%
\opage{366}processes of development, crises and catastrophes in the supersensible world, depicted with such life and invention that when we come to the real, human Fall, it strikes us as being a very insignificant story, a storm in a glass of water. But in that world the author is as in his element. His talent is pronouncedly dogmatic-theosophical, is a talent for thinking in images, a combination of feeling, imagination and idea, better suited to a struggle with language to express the unthinkable and unsayable than to penetrating sympathy with and understanding of real, sharply determined conditions of reality.

It is the great problem of the origin of evil that leads Böhme and with him Martensen up into that supersensible world. In seeing the struggle between good and evil in nature Böhme was seized by deep melancholy and fell into doubt and brooding. How can this struggle be reconciled with God's goodness and unity? Böhme found the solution in the thought, already voiced by Master Eckhart, that without opposition no self-consciousness is possible, or, as Böhme expresses himself in his pithy language: ``Kein Ding ohne Widerwertigkeit mag ihm selber offenbar werden.'' If God is then to be the clearly conscious and willing God, there must be an opposition to overcome, a darkness that can form the background of the light, an agony that can let joy reveal itself, a No that can bring out the Yes. In this ground of nature within God he finds ``the foundation of Hell,'' the germ of evil. By Lucifer's fall it was released and kindled into a fire that burns in everything that is harmful and evil in nature and in man. The whole of this process of development is depicted by the author following Böhme in a grand and captivating manner. One reads the presentation like a poem, a Protestant counterpart to Dante's Divine Comedy (only that here it is a tragedy).

Böhme is a bold thinker, despite all his childlike piety and his Lutheran orthodoxy. He does not shrink back from any thought that lies on his way. Even the biblical presentation he rejects when it ``wider die Philosophia und Vernunft laufet.'' The Mosaic story of creation is adapted according to the Paracelsian philosophy of nature, and the germ of evil is traced back to the innermost nature of the Godhead. Here Martensen lodges his protest. Böhme is to him too naturalistic and pantheistic; Böhme thinks ``after a pagan, pantheistic, cosmogonic type.'' Martensen demands that God be thought as \emph{actual from the first}, actuality prior to all possibility. I cannot see otherwise than that this demand makes all theosophy impossible. Theosophy asks there where most have long since ceased to ask; the actuality of the Godhead presupposes for theosophy, just as all other actuality does, a possibility preceding it (in principle, if not in time). We
% (p. 366, column 2)
\setcounter{footnote}{0}%
understand actuality only through the possibilities from which it unfolds. If the concept of God is to be thought, then, possibility cannot be kept out. Possibility and actuality are concepts of the kind that cannot be separated and used each by itself. It is just as arbitrary to begin with an absolute actuality as with an absolute possibility. We know actuality only in relation to possibility, and conversely. Böhme is right that if we \emph{will} think, we must use the concept of possibility, posit an evolution from possibility to actuality. But he is wrong that we can arrive at any conclusion here: for just as every actuality presupposes a possibility, every possibility presupposes a given actuality. --- Martensen calls Böhme's doctrine pagan. Böhme anticipated such an attack in the following statement (which I have not found in the present book, though it is characteristic): ``Nun wirst du sagen, ich schreibe heidnisch. Höre und siehe und merke den Unterschied wie dies alles sei; denn \emph{ich schreibe nicht heidnisch sondern philosophisch}.'' --- Historically seen, the accusation of paganism is surely incorrect. For it is precisely pagan thinkers (Plato and Aristotle) who have most energetically tried to think the Godhead as eternal actuality; and in Christianity's doctrine of a suffering God the opposite thought emerges.

To avoid all pantheism Martensen maintains that God did not create the world because it was a necessity for him. Precisely for that reason he lays so great weight on the speculations about a supersensible nature, so that the creation of the sensible world can stand as something as arbitrary as possible. It seems as if he here constantly works against himself; for if love is really the motive of creation, then this too must have issued from a need. The love that can do without its work and is equally glad whether it comes to anything or not, we do not think much of here in this ``sensible'' world. Perhaps the higher logic that theosophy commands corrects our ideas of love (as of actuality). Those who are glad to be able tolerably to think according to the ``finite,'' lower logic must then renounce understanding those high things.

Böhme's lasting significance rests not least on the psychological acuteness with which he discovered and held fast the significance of opposition, of contrast, of relativity in all experience, in all mental life and consciousness. His system is a brilliant psychology disguised as theology. It furnishes an indirect proof that our concepts are not suited to the solution of supernatural riddles. Theosophy must first lay the foundation of a new psychology and a new logic before it can reach its goal.
\signature{\textbf{\sffamily Harald Høffding.}}

\end{document}
